\chapter{Literature Review}
\label{ch:lit_rev}

This chapter reviews the literature relevant to the design and motivation of this project. It establishes the foundations of reinforcement learning, discusses model-based approaches, and traces the development of world models leading to DreamerV3~\citep{hafner2023dreamerv3}. It also details the symlog transformation and the specific motivation for this work.

% ================================================================
\section{Reinforcement Learning Foundations}
\label{sec:lit_rl_foundations}
% ================================================================

Reinforcement learning (RL) is a computational framework in which an agent learns to make sequential decisions by interacting with an environment to maximize a cumulative reward signal. The interaction is formalised as a Markov Decision Process (MDP), defined by the tuple $(\mathcal{S}, \mathcal{A}, P, R, \gamma)$, where $\mathcal{S}$ is the state space, $\mathcal{A}$ is the action space, $P(s' \mid s, a)$ is the transition function, $R(s, a)$ is the reward function, and $\gamma \in [0, 1)$ is the discount factor. At each time step $t$, the agent observes a state $s_t \in \mathcal{S}$, selects an action $a_t \in \mathcal{A}$ according to its policy $\pi(a \mid s)$, receives a scalar reward $r_t = R(s_t, a_t)$, and transitions to a successor state $s_{t+1} \sim P(\cdot \mid s_t, a_t)$. The agent's objective is to find a policy $\pi^*$ that maximises the expected discounted return:
\begin{equation}
    G_t = \sum_{k=0}^{\infty} \gamma^k \, r_{t+k}.
    \label{eq:return}
\end{equation}

Central to RL algorithms are value functions. The state-value function $V^\pi(s) = \mathbb{E}_\pi [G_t \mid s_t = s]$ measures the expected return from state $s$ under policy $\pi$. Temporal-difference (TD) learning updates value estimates at each step using the observed reward and the estimated value of the successor state. The standard TD(0) update rule is:
\begin{equation}
    V(s_t) \leftarrow V(s_t) + \alpha \bigl[ r_t + \gamma V(s_{t+1}) - V(s_t) \bigr],
    \label{eq:td_update}
\end{equation}
where $\alpha$ is the learning rate and the quantity in brackets is the TD error $\delta_t$.

In actor-critic architectures, the agent maintains a parameterized policy $\pi_\theta$ (the actor) and a learned value function $V_\phi(s)$ (the critic) that serves as a baseline to reduce gradient variance. The actor parameters are updated using policy gradients:
\begin{equation}
    \nabla_\theta J(\theta) = \mathbb{E}_{\pi_\theta} \Bigl[
        \nabla_\theta \log \pi_\theta(a_t \mid s_t) \, \hat{A}_t
    \Bigr],
    \label{eq:policy_gradient}
\end{equation}
where $\hat{A}_t$ is an estimator of the advantage function $A^\pi(s,a) = Q^\pi(s,a) - V^\pi(s)$.

% ================================================================
\section{Model-Based Reinforcement Learning and World Models}
\label{sec:lit_model_based}
% ================================================================

Model-free RL methods learn policies or value functions directly from observed transitions. Although successful, they typically require large numbers of environment interactions, making them sample-inefficient. Model-based RL addresses this by learning a predictive model of the environment (approximating the transition function $P(s' \mid s, a)$ and the reward function $R(s, a)$) and using it to generate simulated experience.

The Dyna architecture combines real-world acting, model learning, and planning. In Dyna, the learned world model generates simulated transitions that are used to update the policy or value function in the same manner as real experience. By interleaving planning steps with real interactions, the agent improves sample efficiency.

However, model-based planning introduces the challenge of compounding model error. When the learned model is imperfect, errors in predicted transitions accumulate over multi-step rollouts, causing simulated trajectories to diverge from reality. Polices optimized on these inaccurate simulations can perform poorly in the real environment.

To address compounding errors, modern algorithms learn world models and perform planning or "imagination" within a learned state space. The Dreamer family of algorithms has developed this approach for general-purpose applications. The third iteration of this family, DreamerV3~\citep{hafner2023dreamerv3}, introduces scale-invariant learning techniques that enable a single set of hyperparameters to solve diverse environments.

% ================================================================
\section{The Symlog Transformation}
\label{sec:lit_symlog}
% ================================================================

A major challenge in general-purpose RL is the wide variation in reward and value scales. Environments with large reward scales produce large value targets, which generate large gradients in the critic's loss function, destabilising training.

The symmetric logarithm (symlog) transformation, introduced in DreamerV3~\citep{hafner2023dreamerv3}, compresses both positive and negative values logarithmically while preserving the sign:
\begin{equation}
    \operatorname{symlog}(x) = \operatorname{sign}(x) \cdot \ln\bigl(|x| + 1\bigr).
    \label{eq:symlog}
\end{equation}
Its inverse, the symexp function, recovers the original scale:
\begin{equation}
    \operatorname{symexp}(x) = \operatorname{sign}(x) \cdot \bigl(e^{|x|} - 1\bigr).
    \label{eq:symexp}
\end{equation}
The symlog function is smooth, monotonic, and behaves approximately linearly near zero, transitioning to logarithmic compression for large magnitudes. This allows small rewards to retain resolution while large rewards are compressed, protecting gradient stability.

In DreamerV3~\citep{hafner2023dreamerv3}, symlog is applied to the world model's reward predictor and to the critic's target return predictions. This keeps gradients stable across different reward scales.

% ================================================================
\section{Identified Gap and Project Motivation}
\label{sec:lit_gap}
% ================================================================

DreamerV3~\citep{hafner2023dreamerv3} integrates multiple components, including a recurrent state-space model (RSSM), discrete latent representations, KL balancing, unimix distributions, symlog predictions, two-hot encoding, and return normalization. Because these components are evaluated together in complex environments, the specific contribution of the symlog transformation remains entangled with other scaling and planning mechanisms.

This project isolates the effect of the symlog transformation in a simplified setting. By using a feedforward world model in the raw state space, we examine how symlog-transformed losses affect training stability and prevent value divergence across environments with contrasting reward scales.

% ================================================================
\section{Summary}
\label{sec:lit_summary}
% ================================================================

This chapter reviewed reinforcement learning, model-based planning, and the scale-invariance provided by the symlog transformation in DreamerV3~\citep{hafner2023dreamerv3}. We identified the need to isolate the effects of symlog in a simplified architecture. Chapter~\ref{ch:method} details the experimental design and methodology.